\section{6. In vivo functional validation}\label{in-vivo-functional-validation}

\subsection{6.1 Experimental paradigm}\label{experimental-paradigm}

In our rodent DBS protocols, stimulation amplitude is titrated in each animal against an observable effect, and the amplitude at which side effects such as muscle twitches appear is a natural upper reference. A preclinical DBS system should therefore be able to deliver amplitudes in the therapeutic range and the larger amplitudes at which side effects appear. DBS of the cingulate cortex is under active investigation for chronic pain that is refractory to conventional medical management \citep{boccard2017,shirvalkar2025}. We hypothesized that high-amplitude cingulate stimulation would produce freezing or prolonged immobility. We therefore delivered high amplitude settings intended to evoke a robust stimulation-associated behavioral phenotype in vivo.

Six mice implanted with cingulate DBS electrodes (Section 3.1), which had also received a unilateral chronic constriction injury (CCI) of the sciatic nerve, were placed in an open-field chamber (30 cm × 30 cm inner dimensions) and allowed to acclimate for 5 minutes before behavioral recording. Position was tracked with Bonsai \citep{lopes2015} using a webcam mounted above the chamber (Figure 6A). Each mouse was recorded once, in a single session of 17.5--20.7 minutes comprising a baseline period, a period of bilateral high-amplitude DBS and a post-stimulation period. DBS was switched on 5.5--6.9 minutes and off 11.7--14.4 minutes after the start of recording. These mice were separate from those used for the impedance measurements in Section 3.2. Stimulation was delivered by FLEX-DBS with firmware v0.5, as in the rest of this paper, at an amplitude setting above 100 µA, 450 µs pulse duration and 130 Hz pulse frequency (the 130 Hz setting is delivered as 142.5 Hz, Table 2). Electrode impedance and delivered current were not measured in these mice; at the electrode loads measured in Section 3.2, the delivered current would be limited by compliance to approximately 81--101 µA (Figure 4B).

\subsection{6.2 Video-based behavioral analysis}\label{video-based-behavioral-analysis}

The position of each mouse in the open-field chamber was tracked throughout the baseline, DBS and post-stimulation periods (Figure 6C). Video was saved at 320 × 240 px (MPEG-4), and the centroid of the animal was tracked online in an 894 × 877 px tracking frame, with each position exported with its acquisition timestamp. Frame timing was irregular (median approximately 31 frames/s; inter-frame intervals 13--58 ms), so speed was computed as displacement divided by the measured interval between frames, and all averages were weighted by time. Pixel-to-centimeter calibration was measured separately for each recording from its own video, as the wall-to-wall span of the arena floor detected in a median background image (27.30--27.73 px/cm across recordings), because the camera position shifted slightly between sessions. Trajectories were preprocessed in the following order: (1) tracking gaps of up to 0.25 s were linearly interpolated, and longer gaps were left missing; (2) frames implying a speed above 100 cm/s, which corresponded to the tracker jumping to the experimenter's hand, were rejected iteratively and replaced by linear interpolation; (3) a 5-frame median filter was applied; and (4) position was smoothed with a Gaussian filter (standard deviation 0.12 s, applied on the frame index at the median frame rate).

The baseline, DBS and post-stimulation periods were separated by brief handling events in which the experimenter's hand entered the chamber to switch stimulation on and off. Handling events were detected automatically in the video, and detections less than 60 s apart were merged into a single event. Because handling transiently increased locomotion, every analysis window began 90 s after the preceding handling event. This applied to all three periods: the baseline window began 90 s after the start of recording, the DBS window 90 s after the stimulation-on event, and the post-stimulation window 90 s after the stimulation-off event. The 90 s buffer was fixed before the second group of animals was recorded and was not adjusted afterward. Baseline windows ended at the stimulation-on event, DBS windows at the stimulation-off event, and post-stimulation windows at the end of recording. Analyzed durations were 241--324 s (baseline), 263--339 s (DBS) and 225--288 s (post-stimulation).

Immobility was defined by classifying each inter-frame interval as mobile or immobile using speed thresholds with hysteresis. A mobile bout began when speed exceeded 2 cm/s and ended when speed fell below 1 cm/s, with minimum bout durations of 0.5 s (mobile) and 1.0 s (immobile). Percent time immobile was the time-weighted fraction of the analysis window not classified as mobile. The thresholds, the 100 cm/s speed ceiling and the smoothing parameters were chosen values rather than values taken from the literature. Turning was quantified after rediscretizing each trajectory into steps of constant 1 cm length, independent of speed or elapsed time \citep{bovet1988}. A new point was recorded each time the animal had moved 1 cm in a straight line from the previous point, and paths were split at untracked gaps so that no step spanned missing data. The turning angle was the change in heading between consecutive steps, wrapped to ±180 degrees. Turning for each period was the mean absolute turning angle, which counts turns in either direction, and required at least 20 steps (20 cm of path); otherwise it was left undefined. Locality was quantified as the radius of gyration, the time-weighted root-mean-square distance of the animal's position from its mean position. It was computed in consecutive 60 s bins from the start of each analysis window, with bins containing less than 30 s of valid tracking discarded, and averaged across bins weighted by tracked time.

Periods were compared with two-tailed paired t-tests using the animal as the unit of analysis (n = 6). Animals without a defined value in a period were excluded only from comparisons involving that period. For immobility, a Wilcoxon signed-rank test was also computed for baseline versus DBS. To estimate the contribution of within-session habituation, immobility during DBS was additionally compared with the final 2 min of the baseline window. Turning and locality were exploratory analyses developed after the immobility result, and their p-values are reported without correction for multiple comparisons. Analyses were performed in Python; turning, locality, statistics and figures were computed with Python 3.13.5, NumPy 2.5.3, pandas 3.0.6, SciPy 1.18.1 and Matplotlib 3.11.2.

\subsection{6.3 Immobility increased during high-amplitude stimulation}\label{immobility-increased-during-high-amplitude-stimulation}

To confirm that the device delivered stimulation sufficient to produce observable motor effects, mice were recorded in an open field before, during and after FLEX-DBS at high amplitude settings in a single session (Figure 6A). Across the six animals, the percentage of time spent immobile was 46.5 ± 4.1 \% during baseline, 75.7 ± 11.6 \% during DBS and 79.7 ± 8.9 \% after stimulation (mean ± SEM; Figure 6B), a 62.7 \% increase in the group mean from baseline to DBS. Immobility was higher during DBS than at baseline (paired t-test, t(5) = −3.26, p = 0.022) and remained higher after stimulation ended (t(5) = −4.95, p = 0.004). Five of six animals showed the increase; the exception, F1R, became less immobile during DBS (29.4 \% to 19.2 \%), and a Wilcoxon signed-rank test of baseline versus DBS gave p = 0.0625. Because all periods were recorded in one session without a sham control, this increase is associated with stimulation but cannot be separated from within-session habituation. When DBS was compared only with the final 2 min of baseline (50.4 ± 7.3 \% immobile), the difference was smaller and not significant (t(5) = −2.34, p = 0.066), and immobility did not return toward baseline levels after stimulation ended.

In exploratory analyses that were not corrected for multiple comparisons, movement during DBS was more tortuous and more spatially confined than at baseline (Figure 6C--E). Turning per 1 cm step rose from 21.4 ± 1.6 to 32.9 ± 5.9 degrees (t(5) = −2.40, p = 0.062), increasing in five of six animals, again with the exception of F1R. The within-minute radius of gyration fell from 11.0 ± 0.7 to 5.4 ± 1.6 cm (t(5) = 4.13, p = 0.009). Neither measure differed significantly between DBS and post-stimulation or between baseline and post-stimulation: for turning, n = 4 because two animals did not move after stimulation, with p = 0.246 and p = 0.406; for locality, n = 6, with p = 0.405 and p = 0.199. Locality partly reflects immobility, because a stationary animal scores near zero. Together, these data show motor changes during FLEX-DBS at high amplitude settings, consistent with effective stimulation through implanted electrodes. They do not establish that stimulation caused the changes or that the changes reversed when stimulation ended.
